% Generated draft: requires mathematical and source review.
\documentclass[11pt]{article}
\usepackage[margin=25mm]{geometry}
\usepackage{fontspec}
\setmainfont{Latin Modern Roman}
\usepackage{amsmath,amssymb,mathtools,unicode-math}
\setmathfont{Latin Modern Math}
\usepackage{enumitem,needspace,xurl,hyperref}
\setlength{\parindent}{0pt}
\setlength{\parskip}{5pt}
\setlength{\emergencystretch}{3em}
\newcommand{\N}{\mathbb{N}}
\newcommand{\Z}{\mathbb{Z}}
\newcommand{\Q}{\mathbb{Q}}
\newcommand{\R}{\mathbb{R}}
\newcommand{\C}{\mathbb{C}}
\newcommand{\F}{\mathbb{F}}
\newcommand{\A}{\mathbb{A}}
\newcommand{\PP}{\mathbb{P}}
\newcommand{\E}{\mathbb{E}}
\newcommand{\eps}{\varepsilon}
\newcommand{\dd}{\,\mathrm{d}}
\begin{document}
\section*{UnsolvedMath ID: 701340}
\textbf{Source problem code:} catalogue entry
\subsection{Definitions and mathematical statement}
All groups considered are finite, and classes of groups are understood to be closed under isomorphism. A \emph{formation} is a nonempty class of finite groups closed under homomorphic images and finite subdirect products. Here a subdirect product of groups \(G_1,\ldots,G_r\), with \(r\geq 1\), is a subgroup of \(G_1\times\cdots\times G_r\) whose projection onto every factor is surjective. A formation \(\mathfrak H\) is \emph{hereditary} if every subgroup of every member of \(\mathfrak H\) also belongs to \(\mathfrak H\). It is \emph{saturated} if, for every finite group \(G\),
\[
G/\Phi(G)\in\mathfrak H\quad\Longrightarrow\quad G\in\mathfrak H,
\]
where \(\Phi(G)\) is the intersection of the maximal proper subgroups of \(G\), with \(\Phi(1)=1\).

Let \(\mathfrak U\) denote the formation of all finite supersoluble groups. Explicitly, a finite group \(G\) is supersoluble if it has a series
\[
1=G_0\trianglelefteq G_1\trianglelefteq\cdots\trianglelefteq G_n=G
\]
in which every \(G_i\) is normal in \(G\) and every factor \(G_i/G_{i-1}\), for \(1\leq i\leq n\), is cyclic of prime order. The trivial group is included by allowing \(n=0\).

A formation \(\mathfrak F\) is \emph{radical in a class \(\mathfrak H\)} if \(\mathfrak F\subseteq\mathfrak H\) and, for every \(G\in\mathfrak H\) and every pair of normal subgroups \(A,B\trianglelefteq G\),
\[
A,B\in\mathfrak F\quad\Longrightarrow\quad AB\in\mathfrak F,
\qquad AB=\{ab:a\in A,\ b\in B\}.
\]
Define
\[
\mathfrak M=\bigl\{\mathfrak H:\mathfrak H\text{ is a saturated hereditary formation of finite groups and }\mathfrak U\text{ is radical in }\mathfrak H\bigr\}.
\]
The question is whether there exists \(\mathfrak H_{\ast}\in\mathfrak M\) such that
\[
\forall\mathfrak H\in\mathfrak M,\qquad \mathfrak H\subseteq\mathfrak H_{\ast}.
\]
Thus the requested element must contain every member of \(\mathfrak M\). No restriction to soluble ambient groups is imposed. The supplied record labels this question open; that label and the supplied edition information have not been independently verified.
\subsection{Short English statement}
Among saturated hereditary formations of finite groups that contain all finite supersoluble groups and in which the product of any two normal supersoluble subgroups is supersoluble, does one formation contain every other such formation?
\subsection{Sources}
\begin{itemize}
\item Supplied local UnsolvedMath registry record, permanent ID \(701340\), exact source problem code \texttt{catalogue entry}, titled \emph{Kourovka Notebook Problem 17.38}. The supplied status is \texttt{open}.
\item \emph{The Kourovka Notebook (Unsolved Problems in Group Theory)}, 17th issue (2010), Problem 17.38. Proposed by A. F. Vasil'ev and L. A. Shemetkov.
\item The supplied background also identifies the 21st issue edition (2026), \texttt{arXiv:1401.0300v48}. The supplied collection description names E. I. Khukhro and V. D. Mazurov as editors. These bibliographic and edition details are supplied source information, not independently verified references.
\end{itemize}
\end{document}
